\section{Reverse exposure and safe samples}\label{sec:safety}

The guards from lower-weight groups serve two purposes. They determine
the paths, and they protect a sampled label against lower-weight
competition in its layer. We relate these purposes by exposing the guard
mask backwards in auxiliary time. A label leaves its nominal flat at a
step having a reasonably large conditional exit probability with constant
probability. A comparison using the test mask then turns the square of
that exit probability into a safety guarantee.

Fix $H,D$ and a listed group $h$ throughout this section. Thus the groups,
their indices, the density maps, and $Y$ are fixed. As elsewhere, the
deterministic $M,w$ are suppressed from conditioning. Define
\begin{equation}\label{eq:competition-set}
 K_{h,b}=\bigcup_{j<h}
       \bigl(T_j\cap(X_j(b)\setminus X_j(b-1))\bigr).
\end{equation}
This set contains all possible eligible lower-weight competitors at
endpoint $b$, and is contained in $F_{h-1}(b)$. It includes test-mask
memberships on sacrificed labels; this only enlarges the comparison set.
Let $\sigma_h$ count the labels $d\in D_h\cap Y$ whose birth $b$ on $X_h$
satisfies
\begin{equation}\label{eq:safe-definition}
 a_h+2\le b\le0,\qquad b\equiv\varepsilon\pmod2,\qquad
 d\notin Q_h\bigl(F_{h-1}(b-2)+K_{h,b}\bigr).
\end{equation}
These are analysis quantities; the online rule does not compute them.
Our objective is Lemma~\ref{lem:safe}: a constant fraction of
$D_h\cap Y$ contributes to $\sigma_h$ in expectation, apart from the
loss $|D_h|/\kappa$. We first expose the guards backwards so that each
step queries only previously untouched bits. A bound on disjoint
positive certificates then locates an exit of substantial conditional
probability. Comparing two transitions at that exit will establish the
exclusion in \eqref{eq:safe-definition}.

\subsection{The reverse computation and its filtration}

The guard mask was drawn as part of the initial seed and may be fully known
to the adversary. Reverse exposure analyzes that mask's product law; it
introduces no new online randomness and does not restrict the adversary's
information.

For $k<0$, a later state $S=(S_j:j\le h)$ and a guard-bit vector $c$,
define a one-step map $\mathcal G_k(c,S)$ as follows:
\begin{equation}\label{eq:reverse-map}
 \begin{aligned}
 A_0&=F_0(k),\qquad N_j=Q_{j,k}(A_{j-1}),\\
 A_j&=\cl\bigl(N_j+\{e\in G_j\cap S_j:c_e=1\}\bigr)
       \quad(1\le j\le h).
 \end{aligned}
\end{equation}
Its output is $(N_j:j\le h)$. The map is increasing in $k$, in each
coordinate of $S$, and in the bits $c$, by monotonicity and extensivity of
the density maps. We may view $c$ as a bit vector on the fixed union of
the groups, even though only its restrictions to $G_j\cap S_j$ are used.

Starting with $X_j(0)=E$, compute $X(k)$ at
$k=-1,-2,\ldots,a_h-1$ using
\[
 X(k)=\mathcal G_k(C,X(k+1)).
\]
This is exactly the defining path recursion. For an actual later state
$S=X(k+1)$, the new $N_j$ is contained in $S_j$ for every choice of $c$.
Indeed, inductively $A_{j-1}\subseteq S_{j-1}\subseteq F_{j-1}(k+1)$,
so $N_j\subseteq Q_{j,k+1}(F_{j-1}(k+1))=S_j$; all the guard additions
forming $A_j$ also belong to $S_j$. For $j=1$, use
$F_0(k)\subseteq F_0(k+1)$ to begin the induction.

After computing $N_j$, inspect only the bits on
\begin{equation}\label{eq:departing-query}
 G_j\cap(S_j\setminus N_j).
\end{equation}
The omitted additions inside $N_j$ do not change its closure. Fix a label
order within each such batch. Its domain is determined before its first
query: $N_j$ uses only previous groups' bits. The groups are disjoint and
their domains shrink at each reverse step, so no coordinate is ever
queried twice.

Let $\mathcal H_b$ be the sigma-field generated by the complete transcript
just before the step from $X(b)$ to $X(b-1)$, including the queried labels,
their answers, and the states obtained from them. Thus
$\mathcal H_b\subseteq\mathcal H_{b-1}$ as reverse time advances.
Conditional on $\mathcal H_b$, the unqueried bits on
\[
 V_b=\bigcup_{j\le h}(G_j\cap X_j(b))
\]
are jointly independent Bernoulli bits of rate $t$. Indeed, the
unqueried coordinates are exactly these domains, and each preceding
query was chosen from its prior transcript. In particular, the
conditional transition law is \eqref{eq:reverse-map} with fresh product
bits. The transcript reveals no test bits, parity, or flat at a smaller
time that has not yet been computed.

For a fixed nonloop $d$, on $\{d\in X_h(b)\}$ define the predictable exit
probability
\begin{equation}\label{eq:exit-probability}
 p_{b-1}=\Prob\bigl(d\notin X_h(b-1)\mid\mathcal H_b\bigr).
\end{equation}
Set $p_{b-1}=0$ when $d\notin X_h(b)$. Before exit, these probabilities
are nondecreasing in the order of the reverse computation. To see this,
evaluate each conditional transition probability with an independent
product mask $\xi$ on the whole union of the groups. Smaller time and
smaller input state give a smaller output under this same $\xi$.
Taking probabilities proves the assertion even though the domains of
the unqueried bits change.

\subsection{Certificates and exit mass}

To control exits at steps with $p_{b-1}<\eta$, we consider bands of
comparable exit probabilities. We will show that each survival within a
band certifies one fixed increasing event using only that step's newly
queried positive bits. The next lemma bounds how many such disjoint
certificates can occur.

An increasing Boolean event is a family $\mathcal A\subseteq2^V$ closed
under taking supersets. A positive certificate for $\mathcal A$ is a set
$I$ with $I\in\mathcal A$: setting the coordinates in $I$ to one already
forces the event, whatever the other bits are.

\begin{lemma}[Disjoint-query certificates]\label{lem:transactions}
Let $V_0$ be finite, let its bits be independent Bernoulli random
variables, and let $\mathcal A\subseteq2^{V_0}$ be increasing with failure
probability $\delta>0$. An adaptive procedure queries each coordinate at
most once and divides its queries into transactions. The next query,
transaction boundaries, and stopping decisions depend only on the
transcript and, if desired, an independent auxiliary seed. Suppose the
number of transactions has a deterministic finite bound. A transaction
may be declared successful only if the positive bits queried within that
transaction alone form a certificate for $\mathcal A$. Then
\begin{equation}\label{eq:certificate-budget}
 \Ex[\text{number of successful transactions}]
       \le\log(1/\delta).
\end{equation}
Here and below $\log$ denotes the natural logarithm.
\end{lemma}
\begin{proof}
Condition on the auxiliary seed, if present. Conditional on every
possible query transcript, all unqueried bits retain their joint original
product law: the transcript specifies values only at the coordinates it
queries. This follows by induction over the successive, predictably
chosen queries.

For a remaining coordinate set $V\subseteq V_0$, let $p(V)$ be the failure
probability of $\mathcal A$ when the bits on $V$ have their original laws
and all other bits are zero. If $I$ is the set of positive bits already
queried in the current transaction, let $p_I(V)$ be the same failure
probability with the bits on $I$ forced to one. Monotonicity gives
\[
 \delta\le p(V)\le1,\qquad
 0\le p_I(V)\le p(V),\qquad
 p(V\setminus\{e\})\ge p(V).
\]
The ratio $Z=p_I(V)/p(V)$ starts each transaction at one.

Suppose the next query is $e\in V$, whose success probability is $q_e$.
Write $V'=V\setminus\{e\}$, and let $Z'$ be the ratio after that query.
The one-coordinate identity
\[
 p_I(V)=(1-q_e)p_I(V')+q_e p_{I\cup\{e\}}(V')
\]
holds conditional on the full current transcript. In particular, the
denominator $p(V')$ is the same for either answer, so the exact expected
decrement and its logarithmic bound are
\begin{align}
 \Ex[Z-Z'\mid\text{current transcript}]
  &=Z\left(1-\frac{p(V)}{p(V')}\right)\notag\\
  &=Z\left(1-e^{-(\log p(V')-\log p(V))}\right)
    \le\log p(V')-\log p(V).
 \label{eq:log-resource}
\end{align}
The last inequality uses $0\le Z\le1$ and $1-e^{-x}\le x$ for $x\ge0$.

At the end of a successful transaction, the certificate forces
$p_I(V)=0$, so $Z=0$. At every other transaction end, $0\le Z\le1$.
Thus its success indicator is at most its realized net decrement
$1-Z_{\rm end}$. Summing \eqref{eq:log-resource} over its queries and
taking conditional expectations bounds its success probability by the
expected increase of $\log p(V)$. Reset $I$ to the empty set for the next
transaction, but keep the depleted $V$. The logarithmic increments then
telescope across all transactions. Their total is at most
$0-\log p(V_0)=\log(1/\delta)$.

This summation also covers adaptive stopping: there are at most $|V_0|$
queries and a bounded number of transactions, so one may pad every
stopped sequence with zero increments and sum a deterministic finite
number of conditional identities. No stopping-time limit is required.
\end{proof}

\begin{lemma}[Exit mass]\label{lem:exit}
Every nonloop $d$ exits its $h$-th nominal path by the end of the reverse
computation, and
\begin{equation}\label{eq:low-exit-hazard}
 \Prob(\text{$d$ exits at a step with }p_{b-1}<\eta\mid H,D)<\tfrac12.
\end{equation}
Consequently,
\begin{equation}\label{eq:squared-exit-mass}
 \sum_{b=a_h}^{0}
  \Ex\bigl[\ind{d\in X_h(b)}p_{b-1}^{2}\mid H,D\bigr]
       \ge\frac{\eta}{2}.
\end{equation}
\end{lemma}
\begin{proof}
Fix $p_0>0$ and consider steps, before exit, whose conditional exit
probability lies in $[p_0,2p_0)$. By monotonicity they form one block,
possibly empty. Its first step is determined by the pre-step transcript,
since the exit probabilities are predictable. Condition on a possible
transcript there, with time $k_0$ and input state $S^0$. On the fresh bits
in $V_0=\bigcup_{j\le h}(G_j\cap S^0_j)$, fix the increasing event
\[
 \mathcal A=
 \{c:d\in(\mathcal G_{k_0}(c,S^0))_h\}.
\]
Its failure probability $\delta$ is at least $p_0$.

Treat the queries within each subsequent step of the block as one
transaction. If that step does not exit, its own positive queried bits
alone certify $\mathcal A$. Here is the reason that earlier positive
answers are unnecessary. Hold this step's actual incoming state $S$
fixed, and let $I$ consist only of its queried positives. Write
$c^I_e=\ind{e\in I}$. Evaluating $\mathcal G_k(c^I,S)$, with every
other guard bit set to zero, gives
the same output as the actual step. Inductively in $j$, the new $N_j$ is
unchanged, all positive guard additions outside $N_j$ were queried and
belong to $I$, and omitted additions inside $N_j$ are already spanned.
Previously queried coordinates are outside the current domains.
Finally $k\le k_0$ and $S\subseteq S^0$ coordinatewise, so
\[
 \mathcal G_k(c^I,S)
       \subseteq\mathcal G_{k_0}(c^I,S^0).
\]
Non-exit in the actual step therefore implies $I\in\mathcal A$.
The incoming state is held fixed in this comparison; no earlier state
is recomputed after previous positive bits are erased.

All transactions consume distinct coordinates. Lemma~\ref{lem:transactions}
bounds the conditional expected number of non-exit steps in this block
by $\log(1/\delta)\le\log(1/p_0)$. There is at most one exit step. Thus,
also without conditioning on the entry transcript, the expected number
of steps in this range is at most $\log(1/p_0)+1$. Summing their
conditional exit probabilities bounds the probability of exit in this
range by
\[
 2p_0\bigl(\log(1/p_0)+1\bigr).
\]
For $p_0=2^{-(\ell+1)}$, sum this bound over $\ell\ge12$ to obtain
\begin{equation}\label{eq:dyadic-exit}
 \sum_{\ell=12}^{\infty}
   2^{-\ell}\bigl((\ell+1)\log2+1\bigr)
   =2^{-11}(14\log2+1)<\tfrac12.
\end{equation}
An exit at a step with conditional probability zero has probability
zero, since there are finitely many steps. Also $X_h(a_h-1)=L$, so every
nonloop exits. This proves \eqref{eq:low-exit-hazard}.

The total exit mass at steps having $p_{b-1}\ge\eta$ is at least $1/2$.
For those steps $p_{b-1}^{2}\ge\eta p_{b-1}$, and
$\ind{d\in X_h(b)}p_{b-1}$ has expectation equal to the probability of
exit at that step. Summing gives \eqref{eq:squared-exit-mass}.
\end{proof}

\subsection{A hybrid transition}

Lemma~\ref{lem:exit} supplies mass for squared exit probabilities.
To use it, we compare the ordinary guard transition with a second
transition that has the same conditional law but incorporates the test
bits of departing labels. The joint exit of the two transitions will
imply protection against the competitors in $K_{h,b}$.

Fix an endpoint $a_h+2\le b\le0$ and a pre-step transcript
$\mathcal H_b$. Perform the ordinary step from $S=X(b)$ to
$S'=X(b-1)$, obtaining its complete transcript $\mathcal H_{b-1}$.
Both transitions start from the same frozen input $S$; the ordinary
output $S'$ determines which source supplies each hybrid bit:
\begin{equation}\label{eq:hybrid-bits}
 \begin{array}{c|c|c}
  \text{coordinates }e & \text{ordinary guard-bit status}
                    & \widetilde c_e\\ \hline
  G_j\cap S'_j & \text{unqueried} & \ind{e\in C}\\[2pt]
  G_j\cap(S_j\setminus S'_j) & \text{queried in this step}
                    & \ind{e\in T}
 \end{array}
\end{equation}
Let $\widetilde N=\mathcal G_{b-1}(\widetilde c,S)$, and write
$\widetilde A_j$ for its intermediate flats in
\eqref{eq:reverse-map}.

Conditional on the \emph{complete} ordinary-step transcript, the
partition in \eqref{eq:hybrid-bits} is fixed. The retained $C$-bits on
$G_j\cap S'_j$ have never been queried, so all of them jointly retain
their independent rate-$t$ laws. No $T$-bit has been queried, and the
entire test mask is independent of the guard mask. Consequently the
whole hybrid vector on the input domains has the fresh product law,
conditional on $\mathcal H_{b-1}$. Its output law is the ordinary
transition law determined by the earlier $\mathcal H_b$.

In particular, on $\{d\in X_h(b)\}$,
$\Prob(d\notin\widetilde N_h\mid\mathcal H_{b-1})=p_{b-1}$.
The tower property therefore gives
\begin{align}
 &\Prob(d\in X_h(b),\ d\notin X_h(b-1),\ d\notin\widetilde N_h
                         \mid\mathcal H_b)\notag\\
 &\quad=\Ex\bigl[\ind{d\in X_h(b)}\ind{d\notin X_h(b-1)}
                         p_{b-1}\mid\mathcal H_b\bigr]
      =\ind{d\in X_h(b)}p_{b-1}^{2}.
 \label{eq:double-exit}
\end{align}
More generally, the ordinary and hybrid outputs are independent with the
same law conditional on the pre-step transcript. This assertion is for
one endpoint at a time. The same mask $T$ is reused at all endpoints,
and no joint independence of the different hybrid experiments is
asserted or needed.

The hybrid also has a deterministic containment property:
\begin{equation}\label{eq:hybrid-containment}
 Q_h\bigl(F_{h-1}(b-2)+K_{h,b}\bigr)\subseteq\widetilde N_h.
\end{equation}
To prove it, induct over $j<h$. If
$\widetilde A_{j-1}\supseteq F_{j-1}(b-2)$, monotonicity gives
\[
 \widetilde N_j
 =Q_{j,b-1}(\widetilde A_{j-1})
 \supseteq Q_{j,b-2}(F_{j-1}(b-2))=X_j(b-2).
\]
The retained guards in \eqref{eq:hybrid-bits} include exactly
$C_j\cap X_j(b-1)$, which are the guards used to form $F_j(b-2)$.
Hence $\widetilde A_j\supseteq F_j(b-2)$. The hybrid additions also
include every $T$-positive in
$G_j\cap(X_j(b)\setminus X_j(b-1))$. Extensivity
preserves the analogous test additions from earlier groups. Starting
with $F_0(b-1)\supseteq F_0(b-2)$, the induction therefore proves
\[
 \widetilde A_{h-1}
 \supseteq\cl\bigl(F_{h-1}(b-2)+K_{h,b}\bigr).
\]
The focal map $Q_{h,b-1}=Q_h$ is enabled, so another application of
monotonicity proves \eqref{eq:hybrid-containment}. This is an inclusion
for every completed assignment of the bits; it does not condition on
the not-yet-exposed flat $F_{h-1}(b-2)$.

An ordinary exit at this step means that $d$ has birth $b$. If it also
exits the hybrid, \eqref{eq:hybrid-containment} gives its safety
exclusion. Thus \eqref{eq:double-exit} implies
\begin{equation}\label{eq:birth-safety-probability}
 \begin{split}
 &\Prob\bigl(b_h(d)=b,\quad
       d\notin Q_h(F_{h-1}(b-2)+K_{h,b})\mid H,D\bigr)\\
 &\hspace{25mm}\ge
       \Ex\bigl[\ind{d\in X_h(b)}p_{b-1}^{2}\mid H,D\bigr]
       \qquad(a_h+2\le b\le0).
 \end{split}
\end{equation}
These bounds may be summed by linearity. For a fixed label the actual
birth events at distinct endpoints are disjoint.

\begin{lemma}[Safe sampled labels]\label{lem:safe}
For every listed group $h$,
\begin{equation}\label{eq:safe-samples}
 \Ex[\sigma_h\mid H,D]
   \ge\frac{\eta}{4}|D_h\cap Y|-\frac{|D_h|}{\kappa}.
\end{equation}
\end{lemma}
\begin{proof}
Apply Lemma~\ref{lem:exit} to each $d\in D_h\cap Y$. By
Lemma~\ref{lem:paths}, the only birth excluded from the range
$a_h+2\le b\le0$ is the baseline birth $a_h$; a birth at $a_h+1$ is
impossible. Since $p^2\le p$, the total discarded contribution is at
most
\[
 \sum_{d\in D_h\cap Y}\Prob(b_h(d)=a_h\mid H,D)
   =|D_h\cap Y\cap Q_h(L)|
   \le\rk(Q_h(L))\le\frac{|D_h|}{\kappa}.
\]
Here $Y$ is independent in the matroid, and the last inequality is the
generator bound of Lemma~\ref{lem:generators}. Summing
\eqref{eq:birth-safety-probability} over labels and valid birth times
therefore gives a pre-parity expected safe count at least
\[
 \frac{\eta}{2}|D_h\cap Y|-\frac{|D_h|}{\kappa}.
\]
The independent fair parity keeps each such label with probability
$1/2$. Consequently the stronger lower bound
$\eta|D_h\cap Y|/4-|D_h|/(2\kappa)$ holds, which implies
\eqref{eq:safe-samples}.
\end{proof}
